\begin{frame}
 \frametitle{発表の背景}
    数値計算において、関数の最大、最小値をより高精度な値で求めたいという要件が発生することがある。
    その時、微分などを用いて探索を行うのは困難な場合がよくおこるあ。
    そこで、微分情報を用いずに関数の極値を効率的に探索する手法があると嬉しい。
\end{frame}
\begin{frame}
    実際、私の卒業研究のRemesアルゴリズムによる最良近似多項式の計算においても、極値探索が必要となりました。
    はじめは放物線近似を使えばある程度うまくいきそうだと考えて調べると、Brent法という非常に有名な手法があることがを発見した。
    ここではなんか研究過程でたどり着いたBrent法というものが面白かったのでそのアルゴリズムとその歴史について紹介します。
\end{frame}